% SPDX-License-Identifier: MIT
% Ein Buch = eine Sprache; Uebersetzungen liegen unter content/<Sprache>.
\documentclass[11pt,a4paper,twoside,open=right,BCOR=7mm,DIV=12]{scrbook}
\providecommand{\SASDLanguage}{de}
\usepackage{style/sasdmanual}
\input{i18n/\SASDLanguage.tex}
\ifstrequal{\SASDLanguage}{de}{\selectlanguage{ngerman}}{\selectlanguage{english}}
\hypersetup{pdftitle={SASD UI Toolkit -- \ManualBookTitle},
 pdfauthor={SASD GmbH},pdfsubject={SASD UI Toolkit Anwenderhandbuch}}
\begin{document}
\frontmatter
\begin{titlepage}
\pagecolor{SASDDark}\color{white}\thispagestyle{empty}
\vspace*{1.8cm}
{\sffamily\bfseries\fontsize{17}{20}\selectfont SASD UI TOOLKIT\par}
\vspace{2.3cm}
{\sffamily\bfseries\fontsize{37}{42}\selectfont \ManualBookTitle\par}
\vspace{.55cm}
{\sffamily\fontsize{18}{25}\selectfont \ManualBookSubtitle\par}
\vspace{1.1cm}
\begin{tikzpicture}
\draw[SASDTurquoise,line width=1.3pt,rounded corners=5pt] (0,0) rectangle (12,4.4);
\fill[SASDDarkTwo,rounded corners=5pt] (0.03,.03) rectangle (11.97,4.37);
\fill[SASDTurquoise] (.07,3.65) rectangle (11.93,4.30);
\node[anchor=west,SASDDark,font=\ttfamily\bfseries] at (.38,3.96) {SASD UI Toolkit --- Terminal / Rendered / Native};
\node[anchor=west,white,font=\ttfamily\large] at (.55,2.85) {\$ sasd\_ui\_terminal\_demo};
\node[anchor=west,SASDSoft,font=\ttfamily] at (.55,2.15) {Name: \_\_\_\_\_\_\_\_\_\_\_\_};
\node[anchor=west,SASDTurquoise,font=\ttfamily] at (.55,1.43) {[ Begruessen ]   [ Beenden ]};
\node[anchor=west,SASDSoft,font=\ttfamily\small] at (.55,.65) {C++20  |  CMake  |  ANSI/VT  |  SDL3 experimental};
\end{tikzpicture}
\vfill
{\sffamily\Large SASD GmbH\par}
{\sffamily\small \ManualEdition\quad\textbullet\quad \ManualDate\par}
\end{titlepage}
\nopagecolor\color{black}
\cleardoublepage
\thispagestyle{plain}
{\Large\bfseries Impressum und Hinweise\par}
\vspace{.6cm}
\noindent\textbf{Herausgeber:} SASD GmbH\\
\textbf{Projekt:} \url{https://github.com/Robin-Goerlach/SASD-UI-Toolkit}\\
\textbf{Stand:} Repository-main, Dokumentationsbasis vom 12. Oktober 2026.\\
\textbf{Status:} Vorab-Handbuch zu einer noch nicht API-stabilen C++20-Bibliothek.
\par\medskip
Dieses Handbuch ist eine eigenstaendige, neu verfasste Dokumentation. Die
didaktische Gliederung (Einstieg, gefuehrtes Beispiel, Vertiefung,
Nachschlageteil) nimmt klassische Borland-Handbuecher als Anregung,
ohne deren Text, Abbildungen oder Beispielprogramme zu uebernehmen.
Borland und Turbo Vision sind Marken bzw. Werke ihrer Rechteinhaber;
es besteht keine Hersteller- oder API-Kompatibilitaetsaussage.
\par\medskip
C++-Quellcode, Tests und akzeptierte Architekturentscheidungen des Repositorys
gehen diesem Handbuch im Konfliktfall vor. Nicht vorhandene Funktionen
werden nicht als einsatzbereit beschrieben.
\par\medskip
\textbf{Lizenz:} MIT; siehe \url{https://github.com/Robin-Goerlach/SASD-UI-Toolkit/blob/main/LICENSE}.
Die Handbuchquellen stehen ebenfalls unter MIT.
\par\medskip
\textbf{Sprachen:} Diese Ausgabe ist deutsch. Satzvorlage, Sprachtexte und
Kapitelquellen sind getrennt. Weitere Ausgaben koennen unter
\filepath{content/<sprache>/} entstehen; die englische Ausgabe ist
noch nicht uebersetzt.
\cleardoublepage
\tableofcontents
\mainmatter
\input{content/\SASDLanguage/01-einfuehrung.tex}
\input{content/\SASDLanguage/02-installation.tex}
\input{content/\SASDLanguage/03-erste-schritte.tex}
\input{content/\SASDLanguage/04-komponenten.tex}
\input{content/\SASDLanguage/05-layout-events.tex}
\input{content/\SASDLanguage/06-befehle-menues.tex}
\input{content/\SASDLanguage/07-datenmodelle.tex}
\input{content/\SASDLanguage/08-backends.tex}
\input{content/\SASDLanguage/09-praxis.tex}
\input{content/\SASDLanguage/10-fehlersuche.tex}
\appendix
\input{content/\SASDLanguage/anhang.tex}
\backmatter
\cleardoublepage
\chapter*{Weitere Dokumentation}
\addcontentsline{toc}{chapter}{Weitere Dokumentation}
Die verbindlichen Schnittstellen stehen unter \filepath{include/sasd/ui/};
Architekturentscheidungen unter \filepath{docs/adr/}; lauffaehige Demos unter
\filepath{examples/}; Tests unter \filepath{tests/}.
Quelltexte: \url{https://github.com/Robin-Goerlach/SASD-UI-Toolkit}.
\end{document}
